\section{A local criterion for temperedness}\label{sec:local-criterion}

We now apply the local-coefficient comparison to a generic Langlands
quotient. It first supplies adjoint regularity for a completion built
from the tempered inducing representation. A completion of the same
semisimple parameter with pure adjoint will then force the inducing
exponent to vanish.

\begin{proposition}\label{prop:generic-regularity}
Let $G$ be split adjoint absolutely simple over $k$. Suppose a generic
irreducible representation $\tau$ is the Langlands quotient of
$I_P^G(\sigma_\nu)$, where $P=MR$ is a proper standard parabolic,
$\sigma$ is tempered, and $\nu$ is in the open positive chamber for $P$.
Choose a tempered completion $(\rho_\sigma,N_M)$ from
Theorem~\ref{thm:tempered-completion}, twist it by $\nu$, and include
it in $\widehat G$. The resulting pair $(r,N_M)$ has
$r\simeq\rho_\tau$ and
\[
 L(s,\Ad\circ(r,N_M))\quad\text{is regular at }s=1.
\]
\end{proposition}

\begin{proof}
Exactness and heredity of Whittaker coinvariants imply that $\sigma$
is generic. The assertion about $r$ follows from normalized-induction
and twist compatibility in Theorem~\ref{thm:parameters}.

Let $A(z)$ be the unnormalized long intertwiner on
$I_P^G(\sigma_z)$. For tempered inducing data, $A(z)$ is holomorphic
at $\nu$ and $\operatorname{im}A(\nu)$ is the Langlands quotient;
see \cite[Lemme~VII.4.1 and Th\'eor\`eme~VII.4.2]{Renard}
and \cite[Section~2.2]{DCHL}. Since this image is generic, exactness of
Whittaker coinvariants shows that the induced map on the two Whittaker
lines is nonzero at $\nu$. Using the algebraic trivializations with
nonvanishing specializations established in Section~\ref{sec:coefficients},
the local coefficient $C_G(z,\sigma)$, the inverse of that scalar,
is therefore holomorphic and nonzero at $\nu$.

Decompose $\widehat{\mathfrak n}=\bigoplus_bV_b$ under the connected
center of $\widehat M$, and denote its actions by $r_b$. For any
algebraic representation of a tempered parameter, the eigenvalues of
Frobenius on the kernel of monodromy have absolute value at most $1$.
Indeed, a highest weight $m\geq0$ of the $\SL_2$ factor contributes
$Q^{-m/2}$ in \eqref{eq:completion}, and the commuting Weil action
has eigenvalues of absolute value $1$. Consequently its local
$L$-factor is regular and nonzero when the argument has positive real part.

Since $b(\nu)>0$, every denominator in the identity of
Proposition~\ref{prop:comparison}, using $(\rho_\sigma,N_M)$, is
regular and nonzero at $z=\nu$. The product is also regular and
nonzero there. Local $L$-factors have no zeros, so none of the numerator
factors
\[
 L\bigl(1-b(\nu),r_b^\vee\circ(\rho_\sigma,N_M)\bigr)
\]
can have a pole.

Finally the decomposition into the Levi algebra and opposite nilradicals is
\[
 \Lie(\widehat G)=\Lie(\widehat M)
                   \oplus\widehat{\mathfrak n}
                   \oplus\widehat{\mathfrak n}^{-}.
\]
The Killing form identifies the last summand with the dual of the second.
On the first summand the central twist is trivial, so its factor is
regular at $1$ by temperedness. The second contributes arguments
$1+b(\nu)>0$. The third contributes the arguments $1-b(\nu)$ just
controlled by the numerator factors. Multiplication of these factors
proves the assertion.
\end{proof}

\begin{theorem}[Pure completion criterion]\label{thm:local-criterion}
Let $G$ be split adjoint absolutely simple over a local function field
$k$, and let $\tau$ be an irreducible generic smooth representation of
$G(k)$. Suppose $\rho_\tau$ admits a Weil--Deligne completion
$(\rho_\tau,N_*)$ whose adjoint representation is pure of weight zero.
Then $\tau$ is tempered.
\end{theorem}

\begin{proof}
Assume that $\tau$ is not tempered. By the Langlands classification it
has the description in Proposition~\ref{prop:generic-regularity}, with
$P$ proper and $\nu$ strictly positive. There is no real central twisting
direction for an adjoint group. Let $(r,N_M)$ be the completion in that
proposition, retaining its realization in the fixed dual Levi.
Transport the assumed completion by conjugacy to a pair $(r,N_*)$
over this same Weil parameter; its adjoint representation remains pure.

By Lemma~\ref{lem:pure-determinant}, the adjoint $L$-factor of
$(r,N_*)$ is regular at $1$, as is that of $(r,N_M)$ by
Proposition~\ref{prop:generic-regularity}.
Proposition~\ref{prop:open-orbit} therefore conjugates $N_M$ to $N_*$
by the centralizer of $r$. In particular,
$\Ad\circ(r,N_M)$ is pure of weight zero.

The Weil image lies in $\widehat M$, and $N_M$ lies in
$\Lie(\widehat M)$. Thus $\widehat{\mathfrak n}$ is a
Weil--Deligne direct summand of the adjoint representation.
Lemma~\ref{lem:pure-determinant} makes this summand pure of weight
zero and gives
\begin{equation}\label{eq:det-pure}
 \left|\det\bigl(r(\Fr)\mid\widehat{\mathfrak n}\bigr)\right|=1.
\end{equation}

We compute the same determinant before and after the twist. On every
central-weight summand $V_b$, the untwisted tempered parameter has
Weil determinant of absolute value $1$. Its $\SL_2$ factor has determinant
$1$, since $\SL_2$ has no nontrivial algebraic characters. Thus the
untwisted Weil--Deligne Frobenius determinant has absolute value $1$
on $V_b$. Twisting by $\nu$ multiplies each eigenvalue there by
$Q^{-b(\nu)}$. It follows that
\begin{equation}\label{eq:det-positive}
 \left|\det\bigl(r(\Fr)\mid\widehat{\mathfrak n}\bigr)\right|
   =Q^{-\sum_b b(\nu)\dim V_b}<1.
\end{equation}
The inequality is strict because $P$ is proper and every $b(\nu)>0$.
This contradicts \eqref{eq:det-pure}.
\end{proof}

The criterion separates the local issue from the global source of purity.
In particular it requires neither a full local Langlands correspondence
nor an identification of the two monodromy operators in advance.
